\section{Case-Study 1: Projekt SafeLend: Automatisierung der Kreditprüfung} \label{sec:case_study_1}

Die Case-Study \enquote{SafeLend} behandelt die Einführung eines KI-basierten
Kredit-Scorings in einem regulierten Bankumfeld.
Entlang der DASC-PM-Phasen werden zentrale Daten-, Modell- und Prozessrisiken
identifiziert, bewertet und über Risk Gates gesteuert.
Ausgangspunkt sind hoher Zeitdruck, strikte Compliance-Anforderungen
und ein interdisziplinäres Projektsetup.
Der detaillierte Unternehmens- und Projekthintergrund (Zielbild, Datenbasis,
Rahmenbedingungen, Organisation, Ressourcen) ist im Anhang dokumentiert.

\paragraph*{Risikobewertung und Schwellenwerte.}
Für das Scoring wird in dieser Fallstudie die Logik
$\text{Score} = \text{Eintrittswahrscheinlichkeit} \times \text{Schadensauswirkung} \times \text{Entdeckbarkeit}$
auf einer Skala von 1 bis 10 je Dimension genutzt.
Die Entdeckbarkeit beschreibt dabei, wie schwer ein Risiko frühzeitig zu erkennen ist
(1 = leicht, 10 = kaum ohne explizite Analyse).
Als pragmatische Schwellenwerte werden gesetzt: Grün unter 100, Gelb 100 bis 300, Rot über 300.
Das Registerformat entspricht dem Artefakt im Anhang; dort ist auch die Status-Historie dokumentiert.

\medskip
\noindent\textit{Hinweis: Das Szenario SafeLend wurde auf Basis eines KI-generierten Projektkontexts entwickelt
und erhebt keinen Anspruch auf vollständige empirische Abbildung.
Die Risikoanalyse dient der exemplarischen Anwendung der entwickelten RM-Strategie
unter plausiblen, hypothetischen Rahmenbedingungen.}

\subsection{Gedankenexperiment und detaillierte Analyse}

\subsubsection{Phase 1: Problem- und Zieldefinition}
\paragraph*{Ausgangslage und Risiko.}
In Phase 1 werden Zielbild, Risikotoleranz und regulatorische Mindestanforderungen festgelegt.
Hauptgefahr ist eine Zielschieflage, bei der Geschwindigkeit priorisiert wird,
ohne Kontrollanforderungen ausreichend zu definieren.
Als Artefakte entstehen ein erstes Zielsystem, eine Risikoakzeptanzmatrix und ein initialer Governance-Entwurf.
Im Kontext von SafeLend konkurrieren dabei mehrere Interessenlagen:
Das Business drängt auf schnelle Markteinführung,
die Compliance-Abteilung fordert vollständige regulatorische Absicherung,
und Data-Science-Rollen tendieren dazu, Modellfragen vor Governance-Fragen zu priorisieren.
Gleichzeitig sind Entscheidungen aus dieser Phase, etwa die Festlegung von Risikoakzeptanzgrenzen
oder des Entscheidungsrahmens für das KI-Modell,
in späteren Phasen kaum ohne hohe Korrekturkosten revidierbar.
Da ein KI-gestütztes Kreditscoring nach dem EU AI Act als Hochrisikosystem eingestuft wird,
sind bereits ab Phase 1 dokumentierbare Governance-Strukturen regulatorisch verpflichtend.

\paragraph*{Bewertung und Begründung.}
Bewertet wird das Hauptrisiko \textit{Zielschieflage durch Zielkonflikte}:
\begin{itemize}
  \item \textbf{Eintritt (5):} Zielkonflikte zwischen Geschwindigkeit, Fachanforderungen und Compliance sind in frühen Phasen typisch und grundsätzlich steuerbar.
  \item \textbf{Auswirkung (6):} Fehlsteuerungen wirken in alle Folgephasen und verursachen hohe Korrekturkosten.
  \item \textbf{Entdeckbarkeit (6):} Zielschieflagen werden häufig erst nach getroffenen Weichenstellungen sichtbar.
\end{itemize}

\begin{table}[!htbp]
\centering
\begin{tabular}{|p{5.5cm}|c|c|c|c|c|}
\hline
\textbf{Risiko} & \textbf{Eintritt} & \textbf{Auswirkung} & \textbf{Entdeck.} & \textbf{Score} & \textbf{Kategorie} \\
\hline
Zielschieflage durch Zielkonflikte & 5 & 6 & 6 & 180 & Gelb \\
Fehlende Governance-Transparenz & 5 & 5 & 5 & 125 & Gelb \\
Unklare Verantwortlichkeiten in der Freigabe & 4 & 6 & 4 & 96 & Grün \\
\hline
\end{tabular}
\caption{Phase 1: FMEA-basierte Priorisierung der Risiken, eigene Darstellung nach \cite{Bundesministerium_2018}}
\label{tab:cs1_phase1_risiken}
\end{table}

\paragraph*{Maßnahmen.}
Gesteuert wird über verbindliche Governance-Kriterien (Freigaben, dokumentierte
Risikoakzeptanz) und einen Eskalationspfad für Zielkonflikte (Tab.~\ref{tab:cs1_phase1_massnahmen}).

\paragraph*{KRI.}
Der Gate-Status hängt an der Vollständigkeit der Governance-Artefakte: offene Freigabepunkte müssen 0 sein,
die Dokumentationsvollständigkeit der Ziel- und Rollenfestlegung muss 100\% erreichen.

\begin{table}[!htbp]
\centering
\begin{tabular}{|p{4cm}|p{8cm}|}
\hline
\textbf{Maßnahme} & \textbf{Details zur Umsetzung} \\
\hline
Governance-Kriterien & Verbindliche Freigabeprozesse und schriftliche Risikoakzeptanz. \\
Eskalationspfad & Management-Prozess für Konflikte zwischen Business Value, Compliance und Modellleistung. \\
\hline
\end{tabular}
\caption{Phase 1: Spezifische Steuerungsmaßnahmen, eigene Darstellung}
\label{tab:cs1_phase1_massnahmen}
\end{table}

\paragraph*{Gate-Entscheidung.}
\textbf{Freigabe unter Auflagen.} Governance-Maßnahmen starten parallel; Risk Owner werden benannt
und die Entscheidungen im Risikoregister dokumentiert.

\subsubsection{Phase 2: Datenverständnis und Datenvorbereitung}
\paragraph*{Ausgangslage und Risiko.}
In Phase 2 werden die historischen Daten analysiert.
Das zentrale Risiko ist eine unzureichende Datenqualität
und eine systematische Verzerrung (Bias).
Dies entspricht dem Risikotyp \textit{Datenrisiko}.
Kritisch ist, dass historische Kreditdaten möglicherweise gesellschaftliche
oder prozessuale Verzerrungen der Vergangenheit widerspiegeln
(z.B. die systematische Benachteiligung bestimmter Postleitzahlengebiete).
Ein Modell, das auf solchen Daten trainiert wird,
würde diese Verzerrung erlernen und fortschreiben.

\paragraph*{Bewertung und Begründung.}
Bewertet wird das Hauptrisiko \textit{Datenqualität und Bias in historischen Daten}:
\begin{itemize}
  \item \textbf{Eintritt (6):} Historische Datenschnitte sind oft nicht repräsentativ für aktuelle und zukünftige Kontexte.
  \item \textbf{Auswirkung (10):} Eine verzerrte Datengrundlage schwächt das Modell fundamental und kann systematische Fehlentscheidungen mit finanziellen und reputativen Schäden erzeugen.
  \item \textbf{Entdeckbarkeit (9):} Technische Qualitätsfehler werden über Great-Expectations-Checks und Verteilungsprüfungen früh sichtbar; verdeckte Proxy-Bias-Muster bleiben ohne gezielte Fairness-Analysen schwerer erkennbar.
\end{itemize}

\begin{table}[!htbp]
\centering
\begin{tabular}{|p{5.5cm}|c|c|c|c|c|}
\hline
\textbf{Risiko} & \textbf{Eintritt} & \textbf{Auswirkung} & \textbf{Entdeck.} & \textbf{Score} & \textbf{Kategorie} \\
\hline
Datenqualität und Bias in historischen Daten & 6 & 10 & 9 & 540 & Rot \\
Unvollständige Datenabdeckung & 4 & 7 & 5 & 140 & Gelb \\
Datenschutz- und Compliance-Lücken & 5 & 7 & 4 & 140 & Gelb \\
\hline
\end{tabular}
\caption{Phase 2: FMEA-basierte Priorisierung der Risiken, eigene Darstellung nach \cite{Bundesministerium_2018}}
\label{tab:cs1_phase2_risiken}
\end{table}

\paragraph*{Maßnahmen.}
Die Steuerung kombiniert Data-Quality-Gates, Feature-Governance und Bias-Prüfungen;
kritische Merkmale werden im Datenfreigabeprozess gemeinsam durch Fachbereich,
Data Science und Compliance bewertet (Tab.~\ref{tab:cs1_phase2_massnahmen}).

\paragraph*{KRI.}
Als KRI dienen die Pflichtfeldvollständigkeit der Trainingsdaten (Soll: mindestens 98\%)
und die Differenz der Ablehnungsquoten zwischen Postleitzahlen-Clustern (Soll: maximal 5\%).
Great-Expectations-Checks und statistische Verteilungsprüfungen operationalisieren die Entdeckbarkeit.

\begin{table}[!htbp]
\centering
\begin{tabular}{|p{4cm}|p{8cm}|}
\hline
\textbf{Maßnahme} & \textbf{Details zur Umsetzung} \\
\hline
Data-Quality-Gates & Automatisierte Prüfung der Qualitätsregeln und Datenabdeckung. \\
Feature-Governance & Herkunft, Bedeutung und Verantwortung je Merkmal dokumentiert. \\
Bias-Monitoring & Prüfung auf Diskriminierung und Ausschluss instabiler Merkmale. \\
\hline
\end{tabular}
\caption{Phase 2: Spezifische Steuerungsmaßnahmen, eigene Darstellung}
\label{tab:cs1_phase2_massnahmen}
\end{table}

\paragraph*{Gate-Entscheidung.}
\textbf{Stopp und Nachsteuerung.} Freigabe erst nach technischer Umsetzung der
Datenqualitätsregeln und Wirksamkeitsnachweis im Risikoregister.
Nach Umsetzung sinkt der Score auf 300, sodass eine Freigabe unter Auflagen möglich wird.
Offenes Risiko aus Phase 1: Score unverändert, weiterhin Gelb.

\subsubsection{Phase 3: Modellierung und Evaluierung}
\paragraph*{Ausgangslage und Risiko.}
In Phase 3 steht die Modellgüte im Zentrum, einschließlich Segmentstabilität und Fehlerkosten.
Neben Overfitting ist insbesondere die spätere Drift-Anfälligkeit relevant.
Die Evaluierung erfolgt nicht nur global, sondern explizit je Kundensegment,
um verdeckte Schwächen in risikorelevanten Teilgruppen zu erkennen.
Neben klassischen Gütemaßen werden Entscheidungsfolgen betrachtet,
zum Beispiel Fehlklassifikationskosten und Auswirkungen auf das Kreditportfolio.
Besonders relevant im Kreditkontext ist dabei die asymmetrische Fehlerstruktur:
Ein False Positive (zahlungsunfähiger Kunde erhält Kredit) erzeugt direkte Ausfallkosten,
ein False Negative (kreditwürdiger Kunde wird abgelehnt) hingegen
Reputations- und Diskriminierungsrisiken.
Da der EU AI Act automatisierte Kreditentscheidungen als hochriskant einstuft,
muss das Modell für regulatorische Prüfer nachvollziehbar und erklärbar sein.
Die Evaluierungsphase ist damit nicht nur technisch, sondern auch regulatorisch ein zentraler Kontrollpunkt.

\paragraph*{Bewertung und Begründung.}
Bewertet wird das Hauptrisiko \textit{Instabile Modellgüte über Segmente}:
\begin{itemize}
  \item \textbf{Eintritt (5):} Bei der Modellselektion können insbesondere über heterogene Kundensegmente instabile Ergebnisse auftreten.
  \item \textbf{Auswirkung (6):} Fehlklassifikationen wirken im Betrieb direkt auf Kreditentscheidungen und Risikoexposition.
  \item \textbf{Entdeckbarkeit (5):} Schwächen in Teilgruppen können hinter globalen Kennzahlen verborgen bleiben.
\end{itemize}

\begin{table}[!htbp]
\centering
\begin{tabular}{|p{5.5cm}|c|c|c|c|c|}
\hline
\textbf{Risiko} & \textbf{Eintritt} & \textbf{Auswirkung} & \textbf{Entdeck.} & \textbf{Score} & \textbf{Kategorie} \\
\hline
Instabile Modellgüte über Segmente & 5 & 6 & 5 & 150 & Gelb \\
Overfitting bei der Modellselektion & 5 & 6 & 4 & 120 & Gelb \\
Geringe Erklärbarkeit kritischer Entscheidungen & 4 & 5 & 4 & 80 & Grün \\
\hline
\end{tabular}
\caption{Phase 3: FMEA-basierte Priorisierung der Risiken, eigene Darstellung nach \cite{Bundesministerium_2018}}
\label{tab:cs1_phase3_risiken}
\end{table}

\paragraph*{Maßnahmen.}
Steuerung über Robustheits-/Sensitivitätsanalysen, Erklärbarkeit (z.B. SHAP/LIME)
und Mindestkriterien je Segment; Grenzfälle werden manuell geprüft (Tab.~\ref{tab:cs1_phase3_massnahmen}).

\paragraph*{KRI.}
Als KRI gelten eine Segment-AUC von mindestens 0.75 und ein maximaler Performance-Gap zwischen Segmenten von 0.05.

\begin{table}[!htbp]
\centering
\begin{tabular}{|p{4cm}|p{8cm}|}
\hline
\textbf{Maßnahme} & \textbf{Details zur Umsetzung} \\
\hline
Sensitivitätsanalyse & Monte-Carlo-Simulation zur Prüfung der Modellstabilität je Segment. \\
Erklärbarkeit & SHAP/LIME-Analysen zur Identifikation riskanter Muster und Proxy-Effekte. \\
Manuelle Fachprüfung & Verpflichtende Prüfung bei Grenzentscheidungen durch Fachexperten. \\
\hline
\end{tabular}
\caption{Phase 3: Spezifische Steuerungsmaßnahmen, eigene Darstellung}
\label{tab:cs1_phase3_massnahmen}
\end{table}

\paragraph*{Gate-Entscheidung.}
\textbf{Freigabe unter Auflagen.} Drift-relevante Merkmale und Schwellenwerte für das
Betriebsmonitoring werden definiert und dokumentiert.
Carry-Forward-Prüfung: Das Governance-Risiko aus Phase 1 wechselt auf Grün (Regeln/Wege erprobt)
und bleibt unter Beobachtung ohne Maßnahmenpflicht;
das Datenrisiko aus Phase 2 wechselt auf Grün (Qualitätsregeln bewährt, keine kritischen Proxy-Effekte)
und wird weiterhin passiv beobachtet.

\subsubsection{Phase 4: Pilotierung und Nutzbarmachung}
\paragraph*{Ausgangslage und Risiko.}
In Phase 4 wird das System integriert und im Pilotbetrieb unter Realbedingungen beobachtet.
Das Risiko verschiebt sich von Entwicklungsfehlern hin zu Betriebs- und Prozessrisiken,
insbesondere dem \textit{Modellrisiko} durch beginnenden Data Drift.
Schwerpunkt ist die Robustheit der Ende-zu-Ende-Kette (Datenfluss, Schnittstellen)
und die erstmalige Messung von Drift-Indikatoren wie dem Population Stability Index (PSI).
Im Pilot werden bewusst Szenarien mit Lastspitzen oder Ausfallsituationen getestet,
um operative Schwachstellen vor dem breiten Rollout zu identifizieren.
Hinzu kommt eine regulatorische Besonderheit: Das Modell muss vor dem Echtbetrieb
eine formale interne Abnahme durch Compliance und Risikomanagement bestehen,
was eine zusätzliche Dokumentationspflicht für alle Piloterkenntnisse erzeugt.
In der Praxis wird häufig ein Parallelbetrieb durchgeführt, bei dem alte und neue
Entscheidungslogik gleichzeitig laufen und die Ergebnisse verglichen werden,
ein sinnvoller, aber ressourcenintensiver Ansatz zur Kalibrierung des neuen Systems.

\paragraph*{Bewertung und Begründung.}
Bewertet wird das Hauptrisiko \textit{Beginnende Drift im Pilotbetrieb}:
\begin{itemize}
  \item \textbf{Eintritt (5):} Erste Drift-Signale sind im Pilot plausibel, aber noch nicht voll geschäftskritisch.
  \item \textbf{Auswirkung (6):} Ein unerkannter Drift skaliert im Vollbetrieb zu erheblichen Fehlentscheidungen.
  \item \textbf{Entdeckbarkeit (5):} Das KRI-Set ist im Pilot noch nicht vollständig kalibriert; frühe Signale können übersehen werden.
\end{itemize}

\begin{table}[!htbp]
\centering
\begin{tabular}{|p{5.5cm}|c|c|c|c|c|}
\hline
\textbf{Risiko} & \textbf{Eintritt} & \textbf{Auswirkung} & \textbf{Entdeck.} & \textbf{Score} & \textbf{Kategorie} \\
\hline
Beginnende Drift im Pilotbetrieb & 5 & 6 & 5 & 150 & Gelb \\
Schnittstellenfehler in der Ende-zu-Ende-Kette & 4 & 5 & 4 & 80 & Grün \\
Unzureichendes Pilot-Monitoring & 4 & 5 & 5 & 100 & Gelb \\
\hline
\end{tabular}
\caption{Phase 4: FMEA-basierte Priorisierung der Risiken, eigene Darstellung nach \cite{Bundesministerium_2018}}
\label{tab:cs1_phase4_risiken}
\end{table}

\paragraph*{Maßnahmen.}
Für den Übergang in den Betrieb werden Fallback/Rollback, KRI-Eskalationslogik und
regelmäßiges Reporting umgesetzt; Verantwortlichkeiten und Reaktionszeiten für Incidents
werden verbindlich festgelegt (Tab.~\ref{tab:cs1_phase4_massnahmen}).

\paragraph*{KRI.}
Im Pilotbetrieb gilt ein Population Stability Index (PSI) unter 0.20 als grün; die Schnittstellenfehlerquote muss unter 1\% bleiben.

\begin{table}[!htbp]
\centering
\begin{tabular}{|p{4cm}|p{8cm}|}
\hline
\textbf{Maßnahme} & \textbf{Details zur Umsetzung} \\
\hline
Fallback-Prozess & Automatischer Wechsel auf manuelle Bearbeitung bei KI-Ausfall. \\
Rollback-Option & Rückkehr zur Vorversion bei schwerwiegenden Fehlern. \\
KRI-Eskalationslogik & Drift-Indikatoren mit definierten Schwellenwerten und Reportingpflicht. \\
\hline
\end{tabular}
\caption{Phase 4: Spezifische Steuerungsmaßnahmen, eigene Darstellung}
\label{tab:cs1_phase4_massnahmen}
\end{table}

\paragraph*{Gate-Entscheidung.}
Das verbliebene offene Risiko aus Phase 3 wird abschließend geprüft:
Die Pilotphase hat keine kritischen Segmentschwächen gezeigt
und die Drift-Indikatoren liegen stabil unterhalb der definierten Schwellenwerte.
Das Modellrisiko wechselt auf Grün und bleibt unter passiver Beobachtung ohne Maßnahmenpflicht.
Die Phase wird freigegeben, der Übergang in den Vollbetrieb ist zulässig.

\subsubsection{Phase 5: Laufende Steuerung im Betrieb}
\paragraph*{Ausgangslage und Risiko.}
Im Regelbetrieb wird das Szenario zugespitzt:
Das etablierte Monitoring schlägt an.
Der KRI für Data Drift (z.B. der Population Stability Index)
überschreitet den definierten Alarmwert.
Damit steigt der Risiko-Score,
wie im Risikoregister vorgesehen,
automatisch wieder in den roten Bereich.
Parallel zeigen Segmentauswertungen,
dass sich die Modellgüte in wirtschaftlich sensiblen Kundengruppen verschlechtert.
Es liegt also nicht nur ein technisches Signal,
sondern auch ein geschäftskritisches Portfoliorisiko vor.

\paragraph*{Bewertung und Begründung.}
Bewertet wird das Hauptrisiko \textit{Data Drift in kritischen Segmenten}:
\begin{itemize}
  \item \textbf{Eintritt (7):} Im laufenden Betrieb wirken reale Marktveränderungen auf die Merkmalsverteilungen, während die Modellgüte in kritischen Segmenten sinkt.
  \item \textbf{Auswirkung (8):} Technische und fachliche Signale sind gleichzeitig negativ; Fehlentscheidungen und Portfoliorisiken werden real.
  \item \textbf{Entdeckbarkeit (7):} Auswirkungen auf Portfolioebene akkumulieren sich zeitverzögert und sind daher schwerer frühzeitig erkennbar.
\end{itemize}

\begin{table}[!htbp]
\centering
\begin{tabular}{|p{5.5cm}|c|c|c|c|c|}
\hline
\textbf{Risiko} & \textbf{Eintritt} & \textbf{Auswirkung} & \textbf{Entdeck.} & \textbf{Score} & \textbf{Kategorie} \\
\hline
Data Drift in kritischen Segmenten & 7 & 8 & 7 & 392 & Rot \\
Portfoliorisiko durch Fehlklassifikationen & 8 & 8 & 7 & 448 & Rot \\
Verzögerte Incident-Reaktion & 5 & 6 & 5 & 150 & Gelb \\
\hline
\end{tabular}
\caption{Phase 5: FMEA-basierte Priorisierung der Risiken, eigene Darstellung nach \cite{Bundesministerium_2018}}
\label{tab:cs1_phase5_risiken}
\end{table}

\paragraph*{Maßnahmen zur Nachsteuerung und Stabilisierung.}
Als Reaktion auf den erkannten Data Drift
werden die im Risikomanagementplan vorab definierten Sofortmaßnahmen ausgelöst.
Dies folgt den Risikominimierungsstrategien
(vgl. \hyperref[sec:methode2]{Methode~2: Risikominimierungsstrategien}).
Zur Ursachenanalyse wird eine erneute Risikobeurteilung durchgeführt
(vgl. \hyperref[sec:methode1]{Methode~1: Risikobeurteilung}),
insbesondere über Segmentauswertungen,
Robustheitstests (Sensitivitätsanalyse/Monte-Carlo-Simulation)
und Erklärbarkeitsanalysen.
Eine temporäre manuelle Nachprüfung kritischer Fälle wird eingeführt,
um Fehlentscheidungen zu verhindern,
bis das System wieder stabil ist.
Parallel wird ein Re-Training des Modells
mit einem aktualisierten Datenschnitt angestoßen.
Für die Übergangszeit wird die Automatisierungsquote gezielt begrenzt.

Zur nachhaltigen Stabilisierung werden die Monitoring-Prozesse angepasst
(vgl. \hyperref[sec:risiko_monitoring]{Methode~3: Risikomonitoring}):
Monitoring-Takt und KRI-Schwellen werden neu kalibriert,
das Reporting-Intervall verkürzt.
Eine formale Re-Freigabe des vollautomatisierten Betriebs
erfolgt erst nach stabilen Indikatoren über einen definierten Zeitraum.

\paragraph*{KRI.}
Als Alarm-KRI gilt ein PSI über 0.25; zusätzlich wird die Reaktionszeit auf Incidents auf maximal 30 Minuten begrenzt.

\medskip
\noindent
Die RM-Strategie ist in diesem Szenario gut anwendbar: Gate-Entscheidungen strukturieren den Verlauf,
und alle Risiken werden per Carry-Forward-Logik konsequent nachverfolgt, ohne dass jede Abweichung
automatisch zum Projektstopp führt.

\subsection{Reflexion der RM-Strategie am Beispiel SafeLend}

Die Anwendung auf SafeLend zeigt, dass die entwickelte RM-Strategie die realen Herausforderungen von Data-Science-Projekten
in regulierten Kontexten adressiert. Gate-Entscheidungen strukturieren den Verlauf,
und alle Risiken werden über die Carry-Forward-Logik konsequent nachverfolgt, ohne dass jede Abweichung
automatisch zum Projektstopp führt. Dies balanciert Kontrolle und pragmatische Fortschrittssteuerung.

Eine solch umfassende Risikosteuerung erfordert zusätzliche Ressourcen für Governance-Prozesse, Monitoring und tiefergehende Analysen.
Der damit verbundene Aufwand ist jedoch gerechtfertigt. Insbesondere in regulierten Kontexten wie SafeLend
überwiegen die Vorteile erhöhter Kontroll- und Nachweissicherheit gegenüber dem organisatorischen Mehraufwand.